% Peekaboo-VL report: 07_ablation. Included from ../peekaboo_report.tex.
\section{Ablation Study}
\label{sec:ablation}
This section changes one part of the system at a time and measures what
it is worth. Not every change could be tested on the full 9B model: the
first hint experiments ran on the small 0.8B model, and the newer training
experiments use the faster packed recipe (Sec.~\ref{sec:packing}) instead
of the 5.2-hour first 9B run, and the last block uses the final recipe. So
each row in Table~\ref{tab:ablation} names its own reference, and only rows
with the same reference should be compared with each other.

\begin{table}[t]
\centering
\caption{All ablations in one place. Each block changes one part of the
system; ``Ref.'' is the reference score (\%) and $\Delta$ the change in
points. FT = fine-tuned, TW = test-weighted dev.}
\label{tab:ablation}
\footnotesize
\setlength{\tabcolsep}{3pt}
\begin{tabular}{@{}p{0.47\columnwidth}p{0.25\columnwidth}rr@{}}
\toprule
Change & Measured on & Ref. & $\Delta$ \\
\midrule
\multicolumn{4}{@{}l}{\textit{Language model} (Sec.~\ref{sec:backbone}); reference Qwen3.5-9B} \\
Cosmos-Reason2-8B & no FT, dev & 79.6 & $-2.7$ \\
Qwen3.8-27B, 4-bit & no FT, dev & 79.6 & $-2.1$ \\
Qwen3.5-4B & no FT, dev & 79.6 & $-2.6$ \\
Cosmos-Reason2-8B & FT packed, TW & 81.6 & $-0.7$ \\
\midrule
\multicolumn{4}{@{}l}{\textit{Camera-motion hint} (Sec.~\ref{sec:cue}); camera-motion questions (88)} \\
Answer words, not image words & 0.8B no FT & 55.7 & $+15.9$ \\
Answer-worded hint vs.\ no hint & 0.8B no FT & 68.2 & $+3.4$ \\
Hint of a wrong clip & 9B FT & 85.2 & $-13.6$ \\
Hint removed & 9B FT & 85.2 & $-4.5$ \\
\midrule
\multicolumn{4}{@{}l}{\textit{Time detail at test time} (Sec.~\ref{sec:timetest}); exact/range TG (159)} \\
32 frames instead of 16 & 9B FT & 67.3 & $-5.0$ \\
Jitter, 3 frame positions & 9B FT & 67.3 & $-4.4$ \\
\midrule
\multicolumn{4}{@{}l}{\textit{Text prior mixed in} (Sec.~\ref{sec:prior})} \\
$\alpha{=}0.75$ fixed in advance & 0.8B no FT, subset & 71.6 & $+6.0$ \\
$\alpha$ by cross-validation & 9B no FT, dev & 79.6 & $+0.6$ \\
$\alpha$ by cross-validation & 9B FT, dev & 84.7 & $+0.4$ \\
\midrule
\multicolumn{4}{@{}l}{\textit{Training recipe} (Sec.~\ref{sec:packing})} \\
8 questions per video (packing) & 9B FT, TW & 82.4 & $-0.8$ \\
Learning rate $2{\times}10^{-4}$ (was $10^{-4}$) & 9B FT packed, TW & 81.6 & $-3.3$ \\
\midrule
\multicolumn{4}{@{}l}{\textit{Object-motion hint} (Sec.~\ref{sec:objmotion})} \\
Object-motion hint added & 9B FT packed, TW & 81.6 & $-0.9$ \\
\midrule
\multicolumn{4}{@{}l}{\textit{On the final recipe} (Sec.~\ref{sec:final-ablation})} \\
Second training pass & final FT, TW & 83.0 & $-0.1$ \\
Caption hint, dev (PriDe) & final FT & 85.0 & $+0.4$ \\
Caption hint, test questions in val & final FT, 392 Q & 80.6 & $-1.8$ \\
64 frames at test (TG, TP, CD) & final FT, 634 Q & 76.7 & $+0.5$ \\
QA-generator model, Qwen3-VL-30B-A3B & no FT, dev & 79.6 & $+0.4$ \\
\bottomrule
\end{tabular}
\end{table}


\subsection{Which language model?}
\label{sec:backbone}
Our method does not depend on one particular model: the hint, the letter
scoring and the LoRA recipe work with any video language model that reads
frames with timestamps. We therefore tried several backbones
(Table~\ref{tab:backbone}).

Without fine-tuning, Qwen3.5-9B is the strongest. The larger Qwen3.8-27B, run
in 4-bit to fit on our GPU, scores lower (77.5\%), and the smaller
Qwen3.5-4B scores 77.0\%. We also tried NVIDIA's
Cosmos-Reason2-8B~\cite{nvidia2025cosmosreason2}, a Qwen3-VL-8B model with
extra training on real-world physical video such as traffic and robots.
Without fine-tuning it scores 76.9\%, 2.7 points below Qwen3.5-9B; it is
weaker on Viewpoint ($-11.5$) and the reasoning tasks, but slightly better
on Trend \& Pattern ($+2.7$).

We then fine-tuned Cosmos-Reason2-8B with exactly the same recipe as a
Qwen3.5-9B run (8 questions per video in one sequence;
Sec.~\ref{sec:packing}). Fine-tuning helps it even more than Qwen (+7.6
test-weighted), but it still ends 0.7 points below Qwen3.5-9B with the same
recipe and 1.5 points below our first unpacked Qwen3.5-9B model. Almost all of the
gap comes from one task, Change Detection ($-8.7$). On Viewpoint, Trend \&
Pattern, Object \& Land Cover and Temporal Grounding it is 1--2 points
\emph{better}, which makes it a promising partner for combining two models.
Qwen3.5-9B remains our main model and is the one behind our test score.

\begin{table}[t]
\centering
\caption{Language models we tried (dev accuracy \%; TW = test-weighted).
``Packed'' runs put 8 questions per video in one training sequence.}
\label{tab:backbone}
\footnotesize
\setlength{\tabcolsep}{4pt}
\begin{tabular}{@{}lrr@{}}
\toprule
Model & Dev & TW \\
\midrule
\multicolumn{3}{@{}l}{\textit{No fine-tuning}} \\
Qwen3.5-4B & 77.0 & -- \\
Qwen3.8-27B (4-bit) & 77.5 & -- \\
Cosmos-Reason2-8B & 76.9 & 73.3 \\
\textbf{Qwen3.5-9B} & 79.6 & -- \\
Qwen3-VL-30B-A3B (FP8; + caption hint) & \textbf{80.0} & 76.8 \\
\midrule
\multicolumn{3}{@{}l}{\textit{Fine-tuned, packed recipe}} \\
Cosmos-Reason2-8B & 83.4 & 80.9 \\
\textbf{Qwen3.5-9B} & \textbf{84.3} & \textbf{81.6} \\
\midrule
\multicolumn{3}{@{}l}{\textit{Fine-tuned, unpacked}} \\
Qwen3.5-9B, first model (16 frames) & 84.7 & 82.4 \\
\textbf{Qwen3.5-9B}, final (32 frames, plain) & \textbf{85.4} & \textbf{83.0} \\
\bottomrule
\end{tabular}
\end{table}


\subsection{Does the wording of the hint matter?}
\label{sec:cue}
We first tested the motion hint with the small Qwen3.5-0.8B model, without
fine-tuning. We used all 191 dev Viewpoint questions, 88 of which ask about
camera motion, and a control set of 100 causal and general questions that
should not depend on motion (Table~\ref{tab:cue}).

The first wording, which described motion in image terms, made things
\emph{worse}. Accuracy on camera-motion questions fell from 68.2\% to
55.7\%. We looked at the 18 questions that this wording turned from right
to wrong. In 12 of them the correct answer was ``moving forward''. But the
hint never said ``forward''. It said ``pans'' in 16 of them and ``zooms
in'' in 14. The model simply matched words from the hint to words in the
options (``Panning from left to right'', ``Zooming in''). The geometry was
right, but the words pointed to the wrong option.

We then wrote the \emph{same} measurements in the answers' own words. This
raised camera-motion accuracy to 71.6\%. Compared with the image wording,
20 questions went from wrong to right and 6 from right to wrong
($p{=}0.009$). Compared with no hint at all, the gain is small and could be
luck (+3.4 points, $p{=}0.69$). The control questions did not change (89 /
88 / 87\%), so the hint does not distract the model on unrelated
questions.

\begin{table}[t]
\centering
\caption{Effect of the hint's wording, Qwen3.5-0.8B without fine-tuning, dev
(\%): no hint, image wording, camera (answer) wording.}
\label{tab:cue}
\footnotesize
\begin{tabular}{@{}lrrr@{}}
\toprule
Question set ($n$) & No hint & Image & Camera \\
\midrule
Camera-motion Viewpoint (88) & 68.2 & 55.7 & \textbf{71.6} \\
Viewpoint-angle (90) & 41.1 & 45.6 & 44.4 \\
All Viewpoint (191) & 56.0 & 52.4 & \textbf{59.2} \\
Control: causal + general (100) & 89.0 & 88.0 & 87.0 \\
\bottomrule
\end{tabular}
\end{table}

\textbf{Does the trained model actually use the hint?} After fine-tuning,
we scored dev again in two altered ways: with the hint removed, and with
the hint of a different, randomly chosen clip (so every clip gets a wrong
hint; Table~\ref{tab:reliance}). The model uses the hint only where it is
relevant. On camera-motion questions, a wrong hint costs 13.6 points (13
questions lost vs.\ 1 gained, $p{=}0.002$), so the model reads what
the hint says. Removing the hint costs 4.5 points ($p{=}0.34$). So after
training, the frames carry most of the information, and the hint adds a
modest amount where it matters. Change Detection, other Viewpoint questions
and all other tasks do not change (within 1 point). Our pre-registered rule
required a 5-point drop when the hint is removed, so we call this result
inconclusive, but with strong evidence that the model reads the hint's
content.

\begin{table}[t]
\centering
\caption{Does the fine-tuned 9B model rely on the hint? (dev, accuracy \%).}
\label{tab:reliance}
\footnotesize
\begin{tabular}{@{}lrrr@{}}
\toprule
Subset ($n$) & Hint & Removed & Wrong clip's \\
\midrule
Camera-motion Viewpoint (88) & \textbf{85.2} & 80.7 & 71.6 \\
Change Detection (231) & 80.5 & 80.5 & 80.5 \\
Other Viewpoint (103) & 82.5 & 83.5 & 83.5 \\
All other tasks (1,656) & 85.4 & 85.5 & -- \\
\bottomrule
\end{tabular}
\end{table}


\subsection{Does more time detail at test time help?}
\label{sec:timetest}
Section~\ref{sec:h6} showed that wrong exact times are near-misses of
about one timestamp gap. So we tried giving the trained model more time
detail at test time, on its
219 dev temporal questions (Table~\ref{tab:h6}). We tried 32 frames instead
of 16. We also tried ``jitter'': averaging the answers from three views
whose frames are taken at 25\%, 50\% and 75\% of each time segment. Both
made exact-time and range questions \emph{worse}: they fell from 67.3\% to
62.3\% (32 frames) and to 62.9\% (jitter, $p{=}0.039$). Only the relative
questions improved, by up to 6.7 points.

Our explanation: the model was trained on 16 centred frames, so it learned
to read time from that exact layout of timestamps. Moving every timestamp
at test time breaks what it learned. The fix must therefore come during
training: train with 32 frames, or with randomly shifted frames, so the
model learns to read any layout. The final model does both (Step~1), and
its Temporal Grounding accuracy rises from 67.6\% to 70.3\%. For that
model, going further to 64 frames at test time no longer hurts, but it
does not help timing questions either (Sec.~\ref{sec:final-ablation}).

\begin{table}[t]
\centering
\caption{More time detail at test time, dev Temporal Grounding
(fine-tuned 9B, trained at 16 frames; accuracy \%).}
\label{tab:h6}
\footnotesize
\begin{tabular}{@{}lrrr@{}}
\toprule
Condition & All & Exact/range & Relative \\
 & (219) & (159) & (60) \\
\midrule
16 frames, centred (as trained) & \textbf{68.0} & \textbf{67.3} & 70.0 \\
32 frames & 66.2 & 62.3 & \textbf{76.7} \\
Jitter, 3 positions & 65.3 & 62.9 & 71.7 \\
\bottomrule
\end{tabular}
\end{table}


\subsection{Does the text prior help a big model?}
\label{sec:prior}
Mixing the blind scorer into the small 0.8B model (no fine-tuning) added 7.7
points on a 220-question dev subset, with $\alpha$ chosen by
cross-validation (64.1$\to$71.8, $p{=}0.0015$). We repeated the test on 250
new questions with $\alpha{=}0.75$ fixed in advance, and the gain held
(71.6$\to$77.6, $p{=}0.008$). On both sets, accuracy is flat for $\alpha$
between 0.5 and 1.5 (Fig.~\ref{fig:alpha}), so the exact value is not
critical.

With the 9B model, the same mix adds only 0.6 points (79.6$\to$80.1,
$p{=}0.33$). The bigger model already knows most of these text patterns.
After fine-tuning the gain is +0.4 (84.7$\to$85.2, $p{=}0.19$), and the
chosen weight falls to $\alpha{=}0.25$ in every fold. Our pre-registered
rule was to use the mix only if it adds at least 1 point with $p{<}0.05$.
It does not, so the submitted system does not use it.

\begin{figure}[t]
\centering
\begin{tikzpicture}
\begin{axis}[
  width=\columnwidth, height=46mm,
  xlabel={Prior weight $\alpha$}, ylabel={Accuracy (\%)},
  xlabel style={font=\footnotesize}, ylabel style={font=\footnotesize},
  ticklabel style={font=\scriptsize},
  xmin=0, xmax=3.2, ymin=62, ymax=82,
  legend style={font=\scriptsize, at={(0.98,0.04)}, anchor=south east, draw=none, fill=none},
  grid=major, grid style={draw=gray!20},
]
\addplot[s2, thick, mark=*, mark size=1.4pt] coordinates {(0,64.1) (0.25,69.5) (0.5,72.7) (0.75,73.6) (1,73.2) (1.5,73.2) (2,73.2) (3,72.3)};
\addplot[s3, thick, mark=square*, mark size=1.4pt] coordinates {(0,71.6) (0.25,73.6) (0.5,76.0) (0.75,77.6) (1,77.2) (1.5,76.0) (2,74.0) (3,72.4)};
\addplot[s1, thick, mark=triangle*, mark size=1.6pt] coordinates {(0,79.55) (0.1,79.84) (0.25,79.98) (0.5,79.84) (0.75,80.13) (1,79.88)};
\legend{0.8B no fine-tuning (first subset), 0.8B no fine-tuning (repeat subset), 9B no fine-tuning (full dev)}
\end{axis}
\end{tikzpicture}
\caption{Accuracy against the weight of the text prior. The gain is large
and stable for the small model and almost disappears at 9B.}
\label{fig:alpha}
\end{figure}


\subsection{Training recipe: packing and learning rate}
\label{sec:packing}
\textbf{One video, many questions per training example.}
Each training example spends about 73\% of its length on video tokens:
roughly 1,920 of 2,000 tokens. Every question about the same clip repeats
the same video tokens. ReVA has about 15 questions per clip, so most of the
compute is spent reading the same frames again and again.

\emph{Question packing} puts the video in once and then asks eight
questions about it, one after another, as a normal chat:
\[
\text{[video] hint } Q_1 \rightarrow A_1,\; Q_2 \rightarrow A_2,\; \ldots,\; Q_8 \rightarrow ?
\]
The model is trained on the answer letter of every question. Because each
question comes after the previous ones, no masking is needed, and the
running memory of the DeltaNet layers is not a problem. We checked that the
first question's scores are the same packed or unpacked (differences within
bf16 rounding). Training becomes 6.4 times faster per question on a small
model. The risk is a mismatch: during training, later questions see other
questions and answers about the same video, which never happens at test
time.

We repeated our main run with packing; the rule written beforehand was
that accuracy must not drop by more than 1 point (test-weighted dev at
least 81.4\%). The packed run reaches 84.3\% on dev and 81.6\% test-weighted, 0.8 points
below the unpacked run (82.4\%) and inside the 1-point margin. Training
took about 70 minutes instead of 5.2\,h, 4.5 times faster. Temporal
Grounding is slightly higher (68.5\% vs.\ 67.6\%). Packing passes, and
because a run now takes about 70 minutes, the packed run is the
reference for the training experiments below. The 0.8-point cost is most
likely a leak: later questions in a pack can see earlier answers about
the same video, which never happens at test time.

\textbf{Learning rate.} Every run so far used a learning rate of
$10^{-4}$, which we never tuned. A recent study that tuned ten LoRA
variants found that, once the learning rate is tuned, plain LoRA is within
one point of all of them~\cite{lee2026lrmatters}; the learning rate
matters more than the variant. We therefore repeated the packed run with
$2{\times}10^{-4}$, changing nothing else. It reached 78.3\% test-weighted,
3.3 points below the packed run at $10^{-4}$ (81.6\%), so we kept
$10^{-4}$. A lower rate ($5{\times}10^{-5}$) was planned but not run.


\subsection{Does an object-motion hint help?}
\label{sec:objmotion}
The ReVA paper names one main failure of video models: confusing the
camera's movement with the movement of objects. Our camera hint tells the
model how the camera moves, but not how the cars move. The
\emph{object-motion hint} does this. For each pair of frames, we
compute dense optical flow (how every pixel moves), fit the ground's
movement with a homography (a transform that also handles tilted views),
and subtract it. Whatever still moves is a moving object. We keep only
small moving blobs that continue in the same direction in the next step,
and describe them by region of the frame and direction, in the questions'
own words:

\smallskip
\noindent\fbox{\parbox{0.95\columnwidth}{\footnotesize\ttfamily
Object motion (camera motion removed): top-right: moving right to left
(0.9-3.9s).}}
\smallskip

\noindent This targets Change Detection and Trend \& Pattern. The hint is
computed for all 1,047 clips: 32 clips have no clearly moving objects and
the median hint is 40 words long. We keep it only if it raises
test-weighted dev by at least 0.5 points (with $p<0.2$), or raises Change
Detection plus Trend \& Pattern by at least 2 points without any task
losing more than 1.5.

\textbf{Result: no gain.} On the plain pass, test-weighted dev falls from
81.6\% to 80.7\% ($-0.9$). With letter debiasing on both runs it is
unchanged (81.0 vs.\ 81.1); each run answers 62 questions the other
misses ($p{=}1.0$). The tasks move in both directions: Structural Layout
$+3.5$ and Geometric Relation $+2.8$, but Object \& Land Cover $-2.8$ and
Change Detection, one of the two target tasks, $-1.3$. Trend \& Pattern
gains only 1.1. Neither rule is met, so the hint is dropped. Two likely
reasons, not yet checked:
after fine-tuning, the frames already show the model how vehicles move,
and on low, tilted views the residual flow still contains parallax from
buildings that the blob filter does not fully remove.



\subsection{What did not help the final model}
\label{sec:final-ablation}
The last block of Table~\ref{tab:ablation} tests four more ideas on the
final recipe. Besides dev, we use a second check that is closer to the
test: the 392 test questions that match a validation question in clip,
question and option set, so their answers are in the official validation
file. The final model answers 316 of them correctly (80.6\%).

\textbf{A second training pass.} Two passes with the learning rate decayed
over both reach 82.96\% test-weighted dev, against 83.04\% for one pass,
and 315 against 316 of the 392 test questions. One pass is enough.

\textbf{A caption hint from the question writer.} ReVA's questions were
written by a large model from video captions. We replayed that captioning
model (Qwen3-VL-30B-A3B) on every clip and gave its caption as a second
hint. Without fine-tuning and without frames, the caption lifts the 9B
model from 64.5\% to 75.9\% on dev (a different clip's caption: 59.3\%),
so the captions carry much of what the questions ask. After fine-tuning,
however, the gain is gone: dev moves by $+0.4$ (85.5\% vs.\ 85.0\%, within
noise), and on the 392 test questions the model with the caption is
\emph{worse} (309 vs.\ 316). We dropped it.

\textbf{More frames at test time.} Scoring the final model with 64 frames
on its three weakest tasks leaves Change Detection unchanged (79.2\%),
moves Temporal Grounding from 70.3\% to 69.9\%, and raises Trend \& Pattern
from 81.0\% to 83.2\%. The last gain rests on 4 questions ($p{=}0.125$), so
we do not count it.

\textbf{The question writer as a second model.} The same 30B model,
without fine-tuning and with both hints, scores 80.0\% on dev, below every
fine-tuned 9B model. Combined with the final model (weights chosen on dev
with folds split by clip), it adds 0.96 test-weighted points on dev
($p{=}0.076$), but only 1 of the 392 test questions. A combination with the
caption-hint model looked better still on dev ($+0.84$, $p{=}0.014$) and
answered one test question \emph{fewer}. On this benchmark, gains of about
one point on dev did not carry over to the test.
